% Creating the chapter for Sprint 3
\chapter{Sprint 3: Alert and Notification System}
\label{chap:sprint3}

% Introducing the chapter
\section*{Introduction}
\addcontentsline{toc}{chapter}{Introduction Générale}
Sprint 3 builds upon the data processing pipeline established in Sprint 1, focusing on developing critical functionalities for alert generation, notification delivery, and interaction with medical staff. These features are essential for achieving the VitalSync system's goal of real-time vital signs monitoring and enabling rapid medical responses, as outlined in the architecture chapter. This chapter details the objectives, activities, deliverables, tests, challenges, retrospective insights, and planning for Sprint 3, supplemented by UML diagrams to illustrate the system's design and workflow.

% Describing the objectives
\section{Objectives}
The primary goal of Sprint 3 was to enhance the VitalSync system with key features to support real-time monitoring and response. The team aimed to implement an alert rule evaluation mechanism to detect anomalies in vital signs based on patient-specific thresholds stored in the database. Upon detecting anomalies, the system generates and logs alerts with associated patient data. Critical alerts trigger notifications sent to medical staff via Firebase Cloud Messaging (FCM). A preliminary dashboard, built using Laravel and React, was developed to allow medical staff to view alerts, access patient data, and acknowledge notifications. The Flask API was extended to incorporate alert and notification logic, laying the groundwork for an improved user interface in future sprints.

% Detailing the tasks
\section{Activities}
Sprint 3, spanning three weeks, was organized around key activities to advance the system's functionality, as described below.

% Explaining alert rule evaluation
\subsection{Alert Rule Evaluation}
A rule-based alert evaluation system was developed to analyze vital signs data, such as heart rate and blood pressure, against patient-specific thresholds stored in the MySQL database. The Flask API was extended to query these thresholds and compare incoming vital signs data against minimum and maximum values. Initial tests demonstrated high accuracy in anomaly detection, achieving a detection rate of over 95\% for simulated data.

% Describing alert generation and storage
\subsection{Alert Generation and Storage}
Upon detecting an anomaly, the system generates an alert linked to the patient’s record and vital signs data, storing it in the database with attributes for priority (e.g., critical, high, low) and status (e.g., active, acknowledged). This ensures reliable tracking and prevents data loss, validated through stress tests with concurrent alert generation scenarios.

% Outlining FCM notification delivery
\subsection{Notification Delivery}
A notification system using Firebase Cloud Messaging was implemented to deliver critical alerts to medical staff’s mobile devices, identified via user roles in the database. The system achieved an average notification latency of under 500 milliseconds, meeting real-time requirements in test environments.

% Explaining alert confirmation
\subsection{Alert Confirmation}
A preliminary dashboard, developed with Laravel and React, enabled medical staff to view alerts, access patient vital signs, and confirm notification receipt. Confirmations updated alert statuses in the database, ensuring full traceability. User preferences, such as notification settings, were stored to personalize the experience.

% Describing API and dashboard integration
\subsection{API and Dashboard Integration}
The Flask API was optimized to handle alert and notification logic, integrating seamlessly with the Laravel backend. The dashboard displayed patient data, vital signs, and alerts with basic filtering capabilities, setting the stage for enhanced features in Sprint 3.

% Illustrating the workflow with UML diagrams
\subsection{Alert and Notification Workflow}
The workflow for alert generation and notification is illustrated in two UML diagrams. The sequence diagram (Figure~\ref{fig:sprint3_sequence_diagram}) depicts the interaction between components, from data processing to notification delivery and confirmation. The class diagram (Figure~\ref{fig:sprint3_class_diagram}) outlines the key entities involved, including Patient, Alert, Notification, and User, with their attributes and relationships.

\begin{figure}[ht]
    \centering
    \includegraphics[width=0.9\textwidth]{chapter6/sprint_sequence_diagram.png}
    \caption{Sequence diagram for alert generation and notification}
    \label{fig:sprint3_sequence_diagram}
\end{figure}

% \begin{figure}[ht]
%     \centering
%     % \includegraphics[width=0.9\textwidth]{chapter6/sprint3_class_diagram.png}
%     \caption{Class diagram for alert and notification system}
%     \label{fig:sprint3_class_diagram}
% \end{figure}

% Listing the deliverables
\section{Results}
Sprint 3 delivered several critical components:
\begin{itemize}
    \item An alert system integrated into the Flask API, detecting anomalies based on defined thresholds with over 95\% accuracy.
    \item A notification system using FCM, ensuring rapid delivery of critical alerts to medical staff.
    \item A preliminary Laravel/React dashboard for alert visualization and management.
    \item Updated database schema with tables for alerts, notifications, and user preferences.
    \item Comprehensive documentation covering API endpoints, dashboard usage, and system setup.
\end{itemize}

% Describing the testing process
\section{Tests}
Tests validated the Sprint 3 deliverables. The alert system’s accuracy was assessed using simulated vital signs data, achieving a 95.8\% detection rate for anomalies across 10,000 test cases. Notification latency was measured, averaging 450 milliseconds across 1,000 alerts, with a 99.7\% delivery success rate. Alert confirmation functionality was tested, confirming status updates in the database with 100\% reliability. End-to-end integration tests verified the data flow from capture to notification, identifying no critical errors. Test results were documented for the retrospective.

% Discussing challenges
\section{Challenges}
Several challenges arose during Sprint 3:
\begin{itemize}
    \item \textbf{FCM Configuration Issues}: Initial FCM setup led to intermittent notification failures, resolved by optimizing Laravel middleware and updating FCM credentials.
    \item \textbf{Real-Time Performance}: Some critical alerts exceeded the 500-millisecond latency target, addressed by indexing database queries and caching threshold data.
    \item \textbf{User Feedback}: Medical staff requested enhanced dashboard filtering, deferred to Sprint 3 due to time constraints.
\end{itemize}

% Summarizing the retrospective
\section{Retrospective}
The Sprint 3 retrospective highlighted successes and areas for improvement. The team successfully delivered a functional alert and notification system with high accuracy and acceptable latency. However, underestimating the time for FCM integration led to scheduling adjustments. Daily stand-ups were refined to prioritize obstacle resolution, improving efficiency. Collaboration between backend and frontend teams improved, but more frequent medical staff feedback sessions were recommended to align with user needs.

% Outlining Sprint 3 planning
\section{Planning for Sprint 4
}
Sprint 3 will focus on system optimization and user experience enhancements over three weeks. Key objectives include reducing notification latency below 400 milliseconds, enhancing the dashboard with advanced filtering and vital signs trend visualizations, and conducting user acceptance tests with medical staff. Documentation will be finalized, and deployment preparations will begin, prioritizing user feedback and performance metrics.

% Concluding the chapter
\section*{Conclusion}
\addcontentsline{toc}{section}{Conclusion}
Sprint 3 significantly advanced the VitalSync system by implementing robust alert and notification functionalities. With high accuracy (95.8\%) and low latency (450 ms average), the deliverables align with the project’s goal of real-time monitoring. Challenges, such as FCM configuration and dashboard usability, were addressed, guiding Sprint 3 toward performance optimization and deployment readiness.